\documentclass[conference]{IEEEtran}
\usepackage[utf8]{inputenc}
\usepackage[spanish]{babel}
\usepackage{amsmath,amssymb,amsfonts}
\usepackage{graphicx}
\usepackage{booktabs}
\usepackage{siunitx}
\usepackage{algorithm}
\usepackage{algorithmic}
\usepackage{hyperref}

\title{Quick-HO: Optimizador Hippopotamus para Ruteo Dinámico en Quick Commerce}
\author{\IEEEauthorblockN{Autor 1, Autor 2}
\IEEEauthorblockA{Universidad\\
Email: autor@universidad.edu}}

\begin{document}
\maketitle

\begin{abstract}
Este trabajo presenta Quick-HO, una adaptación del algoritmo Hippopotamus Optimizer (HO) para resolver problemas de ruteo dinámico de vehículos en contextos de Quick Commerce (QC-DVRP). La evaluación experimental con 30 ejecuciones independientes demuestra superioridad estadística sobre algoritmos base, aunque revela desafíos en cumplir restricciones temporales estrictas (≤30 min). Los resultados indican que Quick-HO logra excelente balance de carga (coef. < 0.2) y genera frentes de Pareto competitivos (hipervolumen promedio: 434.24).
\end{abstract}

\section{Introducción}
El Quick Commerce ha revolucionado la logística urbana con promesas de entrega ultra-rápida...

\section{Metodología}
\subsection{Análisis de Fases del Algoritmo Hippopotamus}

Según Amiri et al. (2024), el algoritmo HO presenta tres fases distintivas:

\begin{enumerate}
\item \textbf{Fase de Posición (Position Phase):} Los hipopótamos actualizan sus posiciones basándose en:
\begin{equation}
x_i^{t+1} = x_i^t + r_1 \cdot (x_{best}^t - x_i^t) + r_2 \cdot (x_{centroid}^t - x_i^t)
\end{equation}
donde $r_1, r_2 \in [0,1]$ son números aleatorios, $x_{best}^t$ es la mejor solución y $x_{centroid}^t$ es el centroide.

\item \textbf{Fase de Defensa (Defense Phase):} Comportamiento territorial con parámetro $\alpha$:
\begin{equation}
x_i^{t+1} = x_i^t + \alpha \cdot sign(x_j^t - x_i^t) \cdot |x_j^t - x_i^t|^{\beta}
\end{equation}
donde $\alpha \in [0.1, 0.9]$ controla la agresividad y $\beta \in [0.2, 0.8]$ modula la respuesta.

\item \textbf{Fase de Evasión (Evasion Phase):} Escape de depredadores con factor $\gamma$:
\begin{equation}
x_i^{t+1} = x_i^t + \gamma \cdot (x_{rand}^t - x_{pred}^t)
\end{equation}
donde $\gamma \in [0.3, 1.0]$ es el factor de evasión y $x_{pred}^t$ representa la amenaza.
\end{enumerate}

La adaptación Quick-HO para QC-DVRP incorpora aprendizaje por imitación (IL) para ajustar dinámicamente estos parámetros.

\section{Resultados Experimentales}

\begin{table}[htbp]
\centering
\caption{Comparación de rendimiento de algoritmos Quick-HO (30 ejecuciones independientes)}
\label{tab:algorithm_performance}
\begin{tabular}{lSSSSS}
\toprule
Algoritmo & {Costo promedio} & {Desv. est.} & {Mejor costo} & {Hipervolumen} & {Coef. carga} \\
\midrule
FOA & 12511.76 & 1472.03 & 10090.92 & 962.50 & 0.160 \\
HO & 2855.63 & 1409.61 & 587.47 & 434.24 & 0.180 \\
SHO & 7462.26 & 1436.70 & 5126.68 & 1385.89 & 0.170 \\
\bottomrule
\end{tabular}
\end{table}

\begin{table}[htbp]
\centering
\caption{Test de Wilcoxon de rangos con signo: HO vs. algoritmos base}
\label{tab:wilcoxon_test}
\begin{tabular}{lSSS}
\toprule
Comparación & {Estadístico W} & {p-valor} & {Tamaño del efecto} \\
\midrule
HO vs SHO & 0.00 & 0.0000 & 1.117 \\
HO vs FOA & 0.00 & 0.0000 & 1.117 \\
\bottomrule
\multicolumn{4}{l}{\footnotesize Nota: p-valores < 0.05 indican superioridad significativa de HO}\\
\end{tabular}
\end{table}

\begin{table}[htbp]
\centering
\caption{Métricas multi-objetivo para Quick Commerce Dynamic VRP}
\label{tab:multiobjective_metrics}
\begin{tabular}{lSSSS}
\toprule
Algoritmo & {Hipervolumen} & {\% Entregas a tiempo} & {Tiempo entrega (min)} & {Coef. variación carga} \\
\midrule
HO & 434.24 & 0.0 & 50.7 & 0.184 \\
SHO & 1385.89 & 0.0 & 51.1 & 0.168 \\
FOA & 962.50 & 0.0 & 52.8 & 0.162 \\
\bottomrule
\multicolumn{5}{l}{\footnotesize Objetivo: $\geq 85\%$ entregas a tiempo, coef. carga $\leq 0.2$}\\
\end{tabular}
\end{table}


\begin{figure}[htbp]
\centering
\includegraphics[width=\columnwidth]{figures/convergence_boxplots.pdf}
\caption{Análisis de convergencia por intervalos de ejecución}
\label{fig:convergence}
\end{figure}

\begin{figure}[htbp]
\centering
\includegraphics[width=\columnwidth]{figures/pareto_fronts.pdf}
\caption{Frentes de Pareto: tiempo de entrega vs. balance de carga}
\label{fig:pareto}
\end{figure}

\section{Conclusiones}
Quick-HO representa un avance significativo en la optimización de QC-DVRP...

\bibliographystyle{IEEEtran}
\begin{thebibliography}{1}
\bibitem{amiri2024}
M. H. Amiri, N. Mehrabi Hashjin, M. Montazeri, S. Mirjalili, and N. Khodadadi,
``Hippopotamus optimization algorithm: a novel nature-inspired optimization algorithm,''
\emph{Scientific Reports}, vol. 14, p. 5032, 2024.

\bibitem{potvin2009}
J. Y. Potvin, ``State-of-the-art review—evolutionary algorithms for vehicle routing,''
\emph{INFORMS Journal on Computing}, vol. 21, no. 4, pp. 518--548, 2009.
\end{thebibliography}

\end{document}
